% =========================================================
% III. PROBLEM FORMULATION
% =========================================================
\section{Phát biểu bài toán}
\label{sec:problem_formulation}

Nghiên cứu xem xét bài toán hỏi đáp pháp luật nhiều bước, trong đó kết luận có
thể phụ thuộc vào nhiều điều kiện, sự kiện trung gian và hệ quả pháp lý. Gọi
câu hỏi đầu vào là $q\in\mathcal{Q}$. Tri thức pháp luật được biểu diễn bằng đồ
thị có hướng $G=(V,E)$ và tập quy tắc liên kết $\mathcal{R}$. Mỗi nút
$v_i\in V$ là một sự kiện pháp lý đã chuẩn hóa; mỗi cạnh tóm tắt một hoặc nhiều
quy tắc điều kiện--hệ quả. Ngữ nghĩa suy luận cấu trúc được xác định bởi tập quy
tắc, không chỉ bởi tính đạt được trên đồ thị.

Với $q$, hệ thống truy xuất các sự kiện liên quan và xây dựng tập đường suy
luận ứng viên
\begin{equation}
\mathcal{P}(q)=\{P_1,P_2,\ldots,P_K\},
\end{equation}
trong đó
\begin{equation}
P_i=v_0\xrightarrow{r_1}v_1\xrightarrow{r_2}\cdots
\xrightarrow{r_h}v_h.
\end{equation}
Đường suy luận chính được chọn bởi
\begin{equation}
P^*=\arg\max_{P_i\in\mathcal{P}(q)}S(P_i,q),
\end{equation}
trong đó $S(P_i,q)$ kết hợp độ phù hợp ngữ nghĩa và tín hiệu cấu trúc.

\subsection{Ngữ nghĩa quy tắc ba giá trị}

Mỗi sự kiện được gắn với một biến cấu trúc
\begin{equation}
X_i\in\{1,0,\bot\},
\end{equation}
trong đó $1$, $0$ và $\bot$ lần lượt biểu diễn \textsc{true},
\textsc{false} và \textsc{unknown}. Giá trị thứ ba biểu diễn bằng chứng chưa
được xác lập hoặc bị xung đột, không phải một giá trị số nằm giữa 0 và 1. Thiết
kế này dựa trên ngữ nghĩa điểm bất động ba giá trị cho logic programming
\cite{kleene1952metamathematics,fitting1985kripke} và các mở rộng đại số tuyến
tính \cite{sato2020threevalued}.

Để không gán một giá trị số cho $\bot$, trạng thái được mã hóa bằng hai kênh hỗ
trợ
\begin{equation}
t_i=\mathbb{I}[X_i=1],\qquad
f_i=\mathbb{I}[X_i=0].
\label{eq:two_channel_state}
\end{equation}
Do đó, $X_i=\bot$ tương ứng với $(t_i,f_i)=(0,0)$. Nếu đồng thời tồn tại hỗ
trợ dương và âm, engine thận trọng quy trạng thái hiệu dụng về $\bot$.

Gọi $\mathbf{B}^{+},\mathbf{B}^{-}\in\{0,1\}^{|\mathcal{R}|\times |V|}$ là
hai ma trận incidence cho các literal trong thân quy tắc lần lượt yêu cầu trạng
thái đúng và sai. Với quy tắc $r$, đặt
$k_r=\sum_i(B^{+}_{ri}+B^{-}_{ri})$. Thân hội của quy tắc chỉ được thỏa mãn khi
mọi literal có đúng trạng thái yêu cầu:
\begin{equation}
b_r(\mathbf{t},\mathbf{f})=
\mathbb{I}\!\left[
\sum_i B^{+}_{ri}t_i+
\sum_i B^{-}_{ri}f_i=k_r
\right].
\label{eq:conjunctive_body}
\end{equation}
Nếu $e_r$ cho biết ít nhất một ngoại lệ đã thỏa mãn, mức kích hoạt của quy tắc
là $a_r=b_r(1-e_r)$. Gọi $\mathbf{H}^{+}$ và $\mathbf{H}^{-}$ là các ma trận
incidence cho phần đầu quy tắc. Hỗ trợ dương cho sự kiện $i$ được tổng hợp bởi
\begin{equation}
d_i^{+}=\mathbb{I}\!\left[\sum_r H^{+}_{ir}a_r\geq 1\right],
\label{eq:disjunctive_support}
\end{equation}
và hỗ trợ âm được tính tương tự. Vì vậy, các literal trong cùng một quy tắc có
quan hệ hội, còn nhiều quy tắc độc lập có cùng hệ quả biểu diễn các cơ chế thay
thế theo quan hệ tuyển. Cụ thể, $(A\land B)\rightarrow Y$ là một quy tắc có hai
literal trong thân; $(A\lor B)\rightarrow Y$ được chuẩn hóa thành hai quy tắc
$A\rightarrow Y$ và $B\rightarrow Y$. Schema hiện tại không hỗ trợ phần đầu
quy tắc dạng tuyển. Cách nhìn incidence giữa thân và đầu quy tắc là biểu diễn ma
trận thưa tương ứng với các mô hình logic programming trong không gian ma
trận/tensor \cite{sakama2021tensor}.

Trạng thái factual là điểm bất động
\begin{equation}
\mathbf{x}^{F}
=
\operatorname{Fix}(\Phi_{\mathcal{R}},\mathbf{x}^{(0)}),
\label{eq:factual_fixed_point}
\end{equation}
trong đó $\Phi_{\mathcal{R}}$ kích hoạt các cơ chế quy tắc. Các ma trận trên là
biểu diễn đại số chính thức; implementation hiện tại đánh giá các quy tắc thưa
tương đương bằng phép suy luận ký hiệu tới điểm bất động, không thực hiện phép
nhân ma trận dày. Engine hỗ trợ nhiều literal trong thân, nhưng dữ liệu legacy
của thực nghiệm hiện tại ánh xạ mỗi bản ghi thành một antecedent dương và một
consequent dương. Do đó, khả năng xử lý conjunction chưa được đánh giá thực
nghiệm trong benchmark này.

\subsection{Graph ablation và structural intervention}

Gọi $M(P^*)=\{v_1,\ldots,v_{h-1}\}$ là tập sự kiện trung gian của đường chính.
Nghiên cứu phân biệt hai phép toán không tương đương. Thứ nhất, phép kiểm tra tô
pô tạo đồ thị cảm sinh
\begin{equation}
G^{-M}=G[V\setminus\{M\}]
\end{equation}
và kiểm tra tính đạt được có giới hạn
\begin{equation}
\rho_{\mathrm{top}}(s,Y\mid M)
=
\mathbb{I}[Y\text{ đạt được từ }s\text{ trong }G^{-M}].
\label{eq:topological_ablation}
\end{equation}
Phép toán này chỉ cho biết có tìm thấy đường đồ thị thay thế trong giới hạn tìm
kiếm hay không. Nó được dùng như diagnostic/fallback về topology, không phải
định nghĩa phản thực của LegalSCM.

Thứ hai, với câu hỏi yêu cầu loại bỏ rõ ràng một mediator, LegalSCM thực hiện
hard intervention
\begin{equation}
do(X_M=0),
\label{eq:problem_do}
\end{equation}
trong đó $0\equiv\mathrm{FALSE}$. Biến $X_M$ và phần còn lại của cấu trúc quy
tắc vẫn được giữ; mọi quy tắc có phần đầu là $X_M$ bị vô hiệu hóa và cơ chế xác
định $X_M$ được thay bằng phép gán hằng $X_M:=0$. Tương đương,
\begin{equation}
\mathcal{R}^{do(X_M=0)}
=
\{r\in\mathcal{R}\mid \operatorname{head}(r)\neq X_M\}.
\end{equation}
Thế giới phản thực được suy luận lại bởi
\begin{equation}
\mathbf{x}^{CF}_{M}
=
\operatorname{Fix}\!\left(
\Phi_{\mathcal{R}^{do(X_M=0)}},
\mathbf{x}^{(0)}\cup\{X_M=0\}
\right).
\label{eq:cf_fixed_point}
\end{equation}
Với biến đích $Y$, $M$ là cần thiết tương đối với mô hình quy phạm đã mã hóa nếu
$x_Y^{F}=1$ và $x_{Y,M}^{CF}=0$; $M$ không cần thiết nếu
$x_{Y,M}^{CF}=1$; kết quả là chưa xác định nếu một trong hai thế giới có trạng
thái unknown. Đây là hard intervention trên mô hình quy tắc đã mã hóa, không có
bước abduction biến ngoại sinh, và không được đồng nhất với thao tác xóa $M$
khỏi $G$.

Cuối cùng, hệ thống dự đoán nhãn phụ thuộc câu hỏi $z\in\mathcal{Z}$, trong đó
\begin{equation}
\begin{aligned}
\mathcal{Z}=\{&\textsc{supported},\textsc{uncertain},\\
&\textsc{reject-direct-claim}\},
\end{aligned}
\end{equation}
và sinh câu trả lời $a$ cùng các căn cứ $C$:
\begin{equation}
f:(q,G,\mathcal{R})\rightarrow(P^*,z,\mathcal{E}^*,a,C),
\end{equation}
trong đó $\mathcal{E}^*$ là tập bằng chứng đã tinh lọc. Miền nhãn quyết định này
khác với miền trạng thái sự kiện $\{1,0,\bot\}$ và các trạng thái nội bộ như
necessary hoặc non-necessary.
